\documentclass[10pt, letterpaper]{article}

% Packages:
\usepackage[
    ignoreheadfoot, % set margins without considering header and footer
    top=2 cm, % seperation between body and page edge from the top
    bottom=1 cm, % seperation between body and page edge from the bottom
    left=2 cm, % seperation between body and page edge from the left
    right=2 cm, % seperation between body and page edge from the right
    footskip=1.0 cm, % seperation between body and footer
    % showframe % for debugging 
]{geometry} % for adjusting page geometry
\usepackage{titlesec} % for customizing section titles
\usepackage{tabularx} % for making tables with fixed width columns
\usepackage{array} % tabularx requires this
\usepackage[dvipsnames]{xcolor} % for coloring text
\definecolor{primaryColor}{RGB}{0, 79, 144} % define primary color
\usepackage{enumitem} % for customizing lists
\usepackage{fontawesome5} % for using icons
\usepackage{amsmath} % for math
\usepackage[
    pdftitle={Anwesha Ghosh's CV},
    pdfauthor={Anwesha Ghosh},
    pdfcreator={LaTeX with RenderCV},
    colorlinks=true,
    urlcolor=primaryColor
]{hyperref} % for links, metadata and bookmarks
\usepackage[pscoord]{eso-pic} % for floating text on the page
\usepackage{calc} % for calculating lengths
\usepackage{bookmark} % for bookmarks
\usepackage{datetime}
\usepackage{lastpage} % for getting the total number of pages
\usepackage{changepage} % for one column entries (adjustwidth environment)
\usepackage{paracol} % for two and three column entries
\usepackage{ifthen} % for conditional statements
\usepackage{needspace} % for avoiding page brake right after the section title
\usepackage{iftex} % check if engine is pdflatex, xetex or luatex

% Ensure that generate pdf is machine readable/ATS parsable:
\ifPDFTeX
    \input{glyphtounicode}
    \pdfgentounicode=1
    % \usepackage[T1]{fontenc} % this breaks sb2nov
    \usepackage[utf8]{inputenc}
    \usepackage{lmodern}
\fi



% Some settings:
\AtBeginEnvironment{adjustwidth}{\partopsep0pt} % remove space before adjustwidth environment
\pagestyle{empty} % no header or footer
\setcounter{secnumdepth}{0} % no section numbering
\setlength{\parindent}{0pt} % no indentation
\setlength{\topskip}{0pt} % no top skip
\setlength{\columnsep}{0cm} % set column seperation
\makeatletter
\let\ps@customFooterStyle\ps@plain % Copy the plain style to customFooterStyle
% \patchcmd{\ps@customFooterStyle}{\thepage}{
%     \color{gray}\textit{\small Vishnu Varadarajan - Page \thepage{} of \pageref*{LastPage}}
% }{}{} % replace number by desired string
% \makeatother
\patchcmd{\ps@customFooterStyle}{\thepage}{}{}{} % replace number by desired string
\pagestyle{customFooterStyle}

\titleformat{\section}{\needspace{4\baselineskip}\bfseries\large}{}{0pt}{}[\vspace{1pt}\titlerule]

\titlespacing{\section}{
    % left space:
    -1pt
}{
    % top space:
    0.3 cm
}{
    % bottom space:
    0.2 cm
} % section title spacing

\renewcommand\labelitemi{$\circ$} % custom bullet points
\newenvironment{highlights}{
    \begin{itemize}[
        topsep=0.10 cm,
        parsep=0.10 cm,
        partopsep=0pt,
        itemsep=0pt,
        leftmargin=0.4 cm + 10pt
    ]
}{
    \end{itemize}
} % new environment for highlights

\newenvironment{highlightsforbulletentries}{
    \begin{itemize}[
        topsep=0.10 cm,
        parsep=0.10 cm,
        partopsep=0pt,
        itemsep=0pt,
        leftmargin=10pt
    ]
}{
    \end{itemize}
} % new environment for highlights for bullet entries


\newenvironment{onecolentry}{
    \begin{adjustwidth}{
        0.2 cm + 0.00001 cm
    }{
        0.2 cm + 0.00001 cm
    }
}{
    \end{adjustwidth}
} % new environment for one column entries

\newenvironment{twocolentry}[2][]{
    \begin{adjustwidth}{0.2 cm + 0.00001 cm}{0.2 cm + 0.00001 cm}
    \def\secondColumn{#2}
    \setcolumnwidth{\fill, 4.5 cm}
    \begin{paracol}{2}
}{
    \switchcolumn \raggedleft \secondColumn
    \end{paracol}
    \end{adjustwidth}
} % new environment for two column entries

\newenvironment{header}{
    \setlength{\topsep}{0pt}\par\kern\topsep\centering\linespread{1.5}
}{
    \par\kern\topsep
} % new environment for the header

\newcommand{\placelastupdatedtext}{% \placetextbox{<horizontal pos>}{<vertical pos>}{<stuff>}
  \AddToShipoutPictureFG*{% Add <stuff> to current page foreground
    \put(
        \LenToUnit{\paperwidth-2 cm-0.2 cm+0.05cm},
        \LenToUnit{\paperheight-1.0 cm}
    ){\vtop{{\null}\makebox[0pt][c]{
        \small\color{gray}\textit{Last updated in \monthname \ \the\year}\hspace{\widthof{Last updated in May 2025}}
    }}}%
  }%
}%

% save the original href command in a new command:
\let\hrefWithoutArrow\href

% new command for external links:
\renewcommand{\href}[2]{\hrefWithoutArrow{#1}{\ifthenelse{\equal{#2}{}}{ }{#2 }\raisebox{.15ex}{\footnotesize \faExternalLink*}}}


\begin{document}
    \newcommand{\AND}{\unskip
        \cleaders\copy\ANDbox\hskip\wd\ANDbox
        \ignorespaces
    }
    \newsavebox\ANDbox
    \sbox\ANDbox{}

    \placelastupdatedtext
    \begin{header}
        \textbf{\fontsize{24 pt}{24 pt}\selectfont Anwesha Ghosh}

        \vspace{0.3 cm}

        \normalsize
        \mbox{{\color{black}\footnotesize\faMapMarker*}\hspace*{0.13cm}Indore/Sonipat}%
        \kern 0.25 cm%
        \AND%
        \kern 0.25 cm%
        \mbox{\hrefWithoutArrow{mailto:anwesha.ghosh_ug2023@ashoka.edu.in}{\color{black}{\footnotesize\faEnvelope[regular]}\hspace*{0.13cm}anwesha.ghosh\_ug2023@ashoka.edu.in}}%
        \kern 0.25 cm%
        \AND%
        \kern 0.25 cm%
        \AND%
        \kern 0.25 cm%
        % \mbox{\hrefWithoutArrow{https://yourwebsite.com/}{\color{black}{\footnotesize\faLink}\hspace*{0.13cm}yourwebsite.com}}%
        % \kern 0.25 cm%
        % \AND%
        % \kern 0.25 cm%
        % \mbox{\hrefWithoutArrow{https://linkedin.com/in/yourusername}{\color{black}{\footnotesize\faLinkedinIn}\hspace*{0.13cm}yourusername}}%
        % \kern 0.25 cm%
        % \AND%
        % \kern 0.25 cm%
        \mbox{\hrefWithoutArrow{https://anwesha-ghosh7.github.io}{\color{black}{\footnotesize\faFolderOpen[regular]}\hspace*{0.13cm}Personal Website}}%
    \end{header}

    \vspace{0.3 cm - 0.3 cm}
 
    \section{Education}



        
        \begin{twocolentry}{
            
            
        \textit{Aug 2023 – Current}}
            \textbf{Ashoka University}

            \textit{B.Sc. (Hons.) with Research in Mathematics (with a minor in Computer Science). Graduating in May 2027.}
        \end{twocolentry}

        \vspace{0.10 cm}
            \begin{highlights}
                \item \textbf{CGPA:} 3.55/4.0 \href{https://drive.google.com/file/d/1fw_4dC1ZK22tBPrDM5tkYH73lAA70gLm/view?usp=sharing}{transcript as of November 2026}
            \end{highlights}


    \vspace{0.2 cm}

        \begin{twocolentry}{
            \textit{June 2011 - May 2023}}
            \textbf{The Shishukunj International School, Indore}
        \end{twocolentry}

        \begin{onecolentry}
            \begin{highlights}
                \item Scored 95\% in Class 12 CBSE Board Examinations.
                \item \textbf{Subjects:} Mathematics, Computer Science, Economics, History, Political Science and English.
            \end{highlights}
        \end{onecolentry}


    
    
    \section{Research Experience}

    \begin{twocolentry}{\textit{June 2026 - Present}}

        \textbf{Undergraduate Thesis: Coarse Median Spaces}

        Thesis Advisor: \href{https://sites.google.com/view/suraj-krishna}{Dr. Suraj Krishna}
        \vspace{0.10 cm}
    \end{twocolentry}

    \vspace{0.10 cm}
    \begin{twocolentry}{
        \textit{June 2025}
    }
        \textbf{Reading Project: Rational Points on Elliptic Curves}

    \end{twocolentry}


    \vspace{0.10 cm}
    \begin{onecolentry}
        \begin{highlights}
            \item Supervisor: Prof. Bibekanada Maji
            \item Read Rational Points on Elliptic curves by Silverman and Tate. Presented bi-weekly.
        \end{highlights}
    \end{onecolentry}


    \vspace{0.10 cm}


    \begin{twocolentry}{
        \textit{March -- June 2025}
    }
        \textbf{Local Dynamics of Attracting Fixed Points Around Meromorphic Maps}

        \href{https://drive.google.com/file/d/1NMlqgXszyMnFgl4M5DiHI6n90JAXx14p/view?usp=sharing}{Link to paper}
        \href{https://sites.google.com/view/sagar-shrivastava/math-apprenticeship/math-apprenticeship-2025?authuser=0}{Link to programme} 
        \vspace{0.10 cm}

        
    \end{twocolentry}


    \vspace{0.10 cm}
    \begin{onecolentry}
        \begin{highlights}
            \item This is an expository article written under the guidance of \href{https://sites.google.com/view/niladri-patra/home/}{Dr. Niladri Patra}, for the Math Apprenticeship Programme 2025 at Ashoka University, jointly with Vishnu Varadarajan.
            \item Surveys the local behavior of meromorphic maps on the Riemann sphere near attracting fixed points. Introduces key complex-analytic tools, defines Fatou and Julia sets, and focuses on attracting periodic orbits via Koenigs linearization and its global extension. Emphasis is placed on rigorous proofs and geometric intuition in complex dynamics.
        \end{highlights}
    \end{onecolentry}

    \vspace{0.1cm}

    \begin{twocolentry}{
        \textit{Oct 2024}
    }
        \textbf{Methods of Numerical Analysis}

        \href{https://drive.google.com/file/d/1BiHd6AFdqqySM4TEfhrTb8nH-MwguX8t/view?usp=drive_link}{Link to paper}
        \href{https://drive.google.com/file/d/1rLLtRZk3Kwqoicncq2fBckvftPEffnBF/view?usp=drive_link}{Link to poster}
        \href{https://gauravbhatnagar.org/summer-2024-ashoka-math-apprenticeship-programme/}{Link to programme} 

        \vspace{0.10 cm}

        
    \end{twocolentry}


    \vspace{0.10 cm}
    \begin{onecolentry}
        \begin{highlights}
            \item Explores key methods in numerical analysis, focusing on root-finding techniques like bisection, false position, fixed-point iteration, Newton-Raphson, and secant, as well as interpolation methods including Lagrange, Newton, and spline interpolation. The project includes developing an independent Python library that implements these methods without external modules, incorporating basic operations like numerical differentiation.
        \end{highlights}
    \end{onecolentry}

        \vspace{0.2 cm}

 
    \section{Teaching Experience}
        \begin{twocolentry}{
        \textit{Ashoka University}    
            
        \textit{Summer 2026}
        
        \textit{Monsoon 2026}}
            \textbf{Teaching Assistant}

            \textit{FC-0306: Quantitative Reasoning and Mathematical Thinking}
    
        \end{twocolentry}
              


        \vspace{0.10 cm}
        \begin{onecolentry}
            \begin{highlights}
                \item \textbf{Summer 2026:} TA for Prof. Aalok Thakkar. Received TA Rating of 4.77/5 for a class of 50 students. 
                \item \textbf{Monsoon 2026 (Ongoing):} TA for Prof. Sandeep Junjea. Graded quizzes and exams, held weekly office hours. 
            \end{highlights}
        \end{onecolentry}


        \vspace{0.2 cm}


    \begin{twocolentry}{
        \textit{Ashoka University}    
            
        \textit{May 2026 - June 2026}
        
        \textit{May 2025 - June 2025}}
        
        \textbf{Teaching Assistant}
            
            \textit{Number Theory, Lodha Genius Program}

            \textit{Experimental Mathematics, Lodha Genius Program}

        \end{twocolentry}

        \vspace{0.10 cm}
        \begin{onecolentry}
            \begin{highlights}
                \item Assisting in teaching high school students in the Number Theory course by Prof. Shanta Laishram (ISI, Delhi) and Experimental Mathematics course by Prof. Sagar Srivastava (TIFR, Ashoka University).
                \item Taking Tutorials and Office Hours to help students understand the concepts and solve their doubts. Graded assignments and tests.
            \end{highlights}
        \end{onecolentry}


        \vspace{0.2 cm}


    \section{Skills}
    Python, C++, C, Machine Learning, MySQL

    \section{Leadership and Outreach}

        \begin{twocolentry}{\textit{January 2026 -- Present}}
            \begin{highlights}
                \item Member, Ashoka University Women in STEM (AUWS)
            \end{highlights}
        \end{twocolentry}

        \begin{twocolentry}{\textit{June 2024 -- Jun 2025}}
            \begin{highlights}
                \item Vice President, Ashoka Women in Computing Society.
            \end{highlights}
        \end{twocolentry}

        \begin{twocolentry}{\textit{June 2024 -- Jun 2025}}
            \begin{highlights}
                \item Treasurer, Ashoka Mathematics Society.
            \end{highlights}
        \end{twocolentry}
        
        \begin{twocolentry}{\textit{October 2023 -- May 2024}}
            \begin{highlights}
                \item Writer, Kalinga Magazine (Ashoka's Satire Magazine).
            \end{highlights}
        \end{twocolentry}

        \begin{twocolentry}{\textit{January 2024 -- May 2024}}
            \begin{highlights}
                \item Head of Events, Red Bricks Club (Ashoka's Poetry Club).
            \end{highlights}
        \end{twocolentry}
        
        
        \begin{twocolentry}{\textit{May 2022 -- May 2023}}
            \begin{highlights}
                \item Student Editor, Students' Council, The Shishukunj International School.
            \end{highlights}
        \end{twocolentry}

        \begin{twocolentry}{\textit{May 2021 -- May 2022}}
            \begin{highlights}
                \item Editor-in-Chief for the monthly school magazine, The Shishukunj Gazette.
            \end{highlights}
        \end{twocolentry}


    

\end{document}